\documentclass[11pt]{article}

\usepackage[T1]{fontenc}
\usepackage{amsmath,amssymb,amsthm,mathtools}
\usepackage[margin=1in]{geometry}

\newtheorem{lemma}{Lemma}
\newtheorem{proposition}[lemma]{Proposition}

\title{Exact Owner Subtraction and Direct Analytic Closure}
\author{}
\date{Analytic ledger replacement, Packet F}

\begin{document}
\maketitle

\begin{abstract}
The executable exact ledger checks a \(9\times7\) truth table when the
accepted lower and staircase regions are subtracted from the
Round--19 residual. This note replaces all sixty-three state checks and
the five additional interface checks by one set-theoretic calculation.
The surviving set is then placed directly inside the closed domain of the
all-frequency retained-remainder theorem.  Every open and closed face is
retained explicitly, and no artificial frequency interface is introduced.
\end{abstract}

\section{Sets and inherited owners}

Let
\[
 0<\rho_*<\rho_0<\rho_c<\rho_{\rm opt},
 \qquad \rho_{\rm opt}:=\frac{39}{50},
\]
and let \(U(\rho)\), \(k_6(\rho)\), and \(k_{11}(\rho)\) be the upper
owner and the two staircase walls from the main proof. Put
\begin{align}
 D_{\rm low}:={}&
 \left\{\rho_*<\rho\leq\rho_0,
       \ (2\rho)^{-1}<K<U(\rho)\right\}\notag\\
 &\mathbin{\cup}
 \left\{\rho_0<\rho<\rho_c,
       \ d<K<U(\rho)\right\},
 \label{eq:D-low}\\
 D_{\rm high}:={}&
 \left\{\rho_c\leq\rho<\rho_{\rm opt},
       \ k_6(\rho)<K<U(\rho)\right\},
 \label{eq:D-high}
\end{align}
where \(d=(2\rho_0)^{-1}\), and define
\begin{equation}
       D_{19}:=D_{\rm low}\mathbin{\dot\cup}D_{\rm high}.
       \label{eq:D19}
\end{equation}
The dot is justified by the disjoint ratio ranges.

The already-proved inclusive owners used in the subtraction are:
\begin{align}
 L_{\rm own}&:=D_{\rm low},
 \label{eq:lower-owner}\\
 S_{\rm own}&:=
 \left\{\rho_c\leq\rho<\rho_{\rm opt},
          \ k_6(\rho)<K\leq k_{11}(\rho)\right\},
 \label{eq:stair-owner}
\end{align}
The lower owner includes the complete \(\rho=\rho_0\) fibre. The
staircase includes its upper wall \(K=k_{11}(\rho)\). The optical owner
from the main proof has already restricted the high component of
\(D_{19}\) to \(\rho<\rho_{\rm opt}\).  The residual also excludes
\(K=U(\rho)\).

\section{The subtraction}

\begin{proposition}[Exact residual identity]
\label{prop:D20}
One has
\begin{equation}
 \boxed{
 D_{19}\setminus
 (L_{\rm own}\cup S_{\rm own})
 =D_{20},}
 \label{eq:D20-identity}
\end{equation}
where
\begin{equation}
 D_{20}:=
 \left\{\rho_c\leq\rho<\rho_{\rm opt},
          \ k_{11}(\rho)<K<U(\rho)\right\}.
 \label{eq:D20}
\end{equation}
\end{proposition}

\begin{proof}
Since \(L_{\rm own}=D_{\rm low}\) and the union in
\eqref{eq:D19} is disjoint,
\begin{equation}
 D_{19}\setminus L_{\rm own}=D_{\rm high}.
 \label{eq:remove-low}
\end{equation}
Inside \(D_{\rm high}\), the lower frequency inequality is strict:
\(k_6(\rho)<K\). Removing the inclusive staircase
\eqref{eq:stair-owner} therefore gives exactly
\begin{equation}
 D_{\rm high}\setminus S_{\rm own}
 =\left\{\rho_c\leq\rho<\rho_{\rm opt},
          \ k_{11}(\rho)<K<U(\rho)\right\}.
 \label{eq:remove-stair}
\end{equation}
There is no missing equality case: \(K=k_{11}(\rho)\) belongs to
\(S_{\rm own}\), while \(K=U(\rho)\) never belonged to \(D_{19}\).
The face \(\rho=\rho_{\rm opt}\) was already removed when the optical
theorem was applied in forming \(D_{19}\), whereas \(\rho=\rho_c\)
remains in \(D_{\rm high}\).  Thus \eqref{eq:remove-stair} is precisely
\eqref{eq:D20}.
\end{proof}

\section{The five delicate interfaces}

For clarity, the five extra assertions formerly checked separately by
the executable ledger are consequences of the displayed sets:
\begin{enumerate}
\item At \(\rho=\rho_0\), the identity
      \((2\rho_0)^{-1}=d\) makes the prospective second lower fibre empty;
      the first set in \eqref{eq:D-low} owns the full equality fibre, and
      \(L_{\rm own}\) removes it exactly once.
\item The face \(\rho=\rho_c\) belongs to \(D_{\rm high}\). It is not an
      optical face, so points with
      \(k_{11}(\rho_c)<K<U(\rho_c)\) remain in \(D_{20}\).
\item The face \(\rho=\rho_{\rm opt}=39/50\) belongs to
      the optical theorem and is already absent from \(D_{19}\), hence
      from \(D_{20}\).
\item The frequency face \(K=k_{11}(\rho)\) belongs to
      \(S_{\rm own}\), whereas \(K=U(\rho)\) is excluded already by the
      strict inequality in \(D_{19}\).
\item If \(\rho_c<\rho<\rho_{\rm opt}\) and
      \(k_{11}(\rho)<K<U(\rho)\), the point is in \(D_{\rm high}\) but in
      neither \(L_{\rm own}\) nor \(S_{\rm own}\); hence it survives,
      and every surviving point has exactly these properties.
\end{enumerate}

\section{Direct retained-remainder ownership}

Let
\begin{equation}
 \mathcal R_{\rm rr}:=
 \left\{(\rho,K):
  \rho_c\leq\rho\leq\rho_{\rm opt},\quad
  K\geq k_{11}(\rho)\right\}.
 \label{eq:R-rr}
\end{equation}
This is the closed domain of the all-frequency retained-remainder theorem
proved by the structural, angular-block, and scalar-positivity modules.
In the notation of the main proof, that theorem states
\begin{equation}
 (\rho,K)\in\mathcal R_{\rm rr}
 \quad\Longrightarrow\quad
 N_D(A_{\rho,1},K^2)<W(\rho,K).
 \label{eq:all-K-rr}
\end{equation}

\begin{proposition}[Direct ownership of the exact residual]
\label{prop:direct-closure}
One has
\begin{equation}
 \boxed{D_{20}\subset\mathcal R_{\rm rr}},
 \qquad
 D_{20}\setminus\mathcal R_{\rm rr}=\varnothing.
 \label{eq:direct-closure}
\end{equation}
Consequently the strict P\'olya inequality holds at every point of
\(D_{20}\).
\end{proposition}

\begin{proof}
If \((\rho,K)\in D_{20}\), then
\[
 \rho_c\leq\rho<\rho_{\rm opt}
 \quad\text{and}\quad
 k_{11}(\rho)<K<U(\rho).
\]
The first pair of inequalities implies
\(\rho_c\leq\rho\leq\rho_{\rm opt}\), and the second implies
\(K\geq k_{11}(\rho)\).  Hence \((\rho,K)\in\mathcal R_{\rm rr}\),
which proves \eqref{eq:direct-closure}; the final assertion follows from
\eqref{eq:all-K-rr}.

The inclusion does not change any inherited owner.  The face
\(\rho=\rho_{\rm opt}\) and the frequency face
\(K=k_{11}(\rho)\) are outside \(D_{20}\), although the stronger theorem
also proves them.  The face \(K=U(\rho)\) remains excluded from
\(D_{20}\), while \(\rho=\rho_c\) remains included subject to the two
strict frequency inequalities.  There is no additional frequency face.
\end{proof}

\paragraph{Ledger rows replaced.}
Proposition~\ref{prop:D20} replaces all \(9\cdot7=63\) state evaluations
in the executable family \texttt{verify\_exact\_subtraction}. The five
items above replace its five separately asserted interface predicates.
Proposition~\ref{prop:direct-closure} supplies the final residual closure
by one set inclusion, without a frequency classifier or executable truth
table.

\end{document}
